\chapter{Phase 3: DOSBox-Staging and Hugo}

\section{Goal}

Have Hugo running on the Pi, displayed on the HDMI screen and playable with a USB keyboard. This is the last individual piece before tying everything together in Phase 4.

\section{Legal status of the game}

Hugo is \emph{abandonware}: commercial software whose support by the original publisher ceased decades ago, and whose passive distribution is broadly tolerated though not formally authorized. Sites like \texttt{myabandonware.com} and \texttt{abandonwaregames.net} host downloadable copies. The legal situation varies by jurisdiction, but home and educational use has not been the subject of litigation.

\section{Step 1: Download Hugo}

From the Mac, open in a browser:

\begin{plaincode}
https://www.myabandonware.com/game/hugo-eya
\end{plaincode}

Download the ZIP file the site provides. Once in \texttt{\~{}/Downloads/}, transfer it to the Pi:

\begin{shellcode}
scp ~/Downloads/Hugo*.zip beto@hugo-pbx.local:/tmp/
\end{shellcode}

On the Pi, unzip into the game's directory:

\begin{shellcode}
sudo mkdir -p /opt/hugo-phone/games/hugo
sudo unzip /tmp/Hugo*.zip -d /opt/hugo-phone/games/hugo
sudo chown -R beto:beto /opt/hugo-phone/games/hugo
ls /opt/hugo-phone/games/hugo/
\end{shellcode}

You should see \texttt{hugo.exe} (or a similar file). Some distributions ship an \texttt{install.exe} installer that must be run inside DOSBox the first time.

\section{Step 2: Verify DOSBox-Staging configuration}

The install script copied a \texttt{dosbox-staging.conf} optimized for Hugo onto the Pi. Verify:

\begin{shellcode}
cat ~/.config/dosbox/dosbox-staging.conf | head -30
\end{shellcode}

If the game's executable is not exactly \texttt{hugo.exe}, edit the \texttt{[autoexec]} section:

\begin{shellcode}
nano ~/.config/dosbox/dosbox-staging.conf
\end{shellcode}

The section to tweak is:

\begin{plaincode}
[autoexec]
mount c /opt/hugo-phone/games
c:
cd hugo
hugo.exe
exit
\end{plaincode}

\section{Step 3: Graphical output modes}

DOSBox-Staging can draw in several ways. In order of increasing simplicity:

\subsection{Development mode: with a lightweight desktop}

For tests and diagnostics, install Openbox (a minimalist window manager):

\begin{shellcode}
sudo apt install -y xinit openbox
echo "exec openbox-session" > ~/.xinitrc
startx /usr/games/dosbox-staging
\end{shellcode}

This starts an X11 server with Openbox and launches DOSBox-Staging inside. Useful for experimentation.

\subsection{Appliance mode: direct KMSDRM}

For the final operation, DOSBox-Staging can draw directly to the kernel framebuffer with no X server, which reduces memory usage and removes compositing latency:

\begin{shellcode}
SDL_VIDEODRIVER=kmsdrm dosbox-staging
\end{shellcode}

\begin{note}
KMSDRM only works from a local TTY (not over SSH X-forwarding). It is the recommended mode for normal appliance operation.
\end{note}

\section{Step 4: Auto-start the game}

Configure the Pi so that powering on takes you directly into Hugo, with no user intervention.

\subsection{Enable autologin on TTY1}

\begin{shellcode}
sudo raspi-config
\end{shellcode}

Navigate to: \textsf{System Options} $\to$ \textsf{Boot / Auto Login} $\to$ \textsf{Console Autologin}. Save and exit.

\subsection{Run DOSBox at login}

Edit the \texttt{.bashrc} of the \texttt{beto} user:

\begin{shellcode}
nano ~/.bashrc
\end{shellcode}

Append at the end:

\begin{plaincode}
if [ -z "$DISPLAY" ] && [ "$(tty)" = "/dev/tty1" ]; then
  SDL_VIDEODRIVER=kmsdrm dosbox-staging
fi
\end{plaincode}

The condition prevents DOSBox from running if you log in over SSH or from a TTY other than 1.

After reboot, the sequence should be:

\begin{enumerate}
  \item Kernel boot
  \item Autologin as \texttt{beto} on \texttt{tty1}
  \item DOSBox-Staging running in KMSDRM mode
  \item Hugo auto-launched per \texttt{[autoexec]}
\end{enumerate}

\section{Step 5: Playability test}

With a USB keyboard plugged into the Pi and Hugo loaded, test:

\begin{itemize}
  \item Arrow keys $\to$ character movement
  \item Space bar $\to$ action
  \item \texttt{Esc} $\to$ menu/exit
\end{itemize}

If Hugo feels excessively fast or slow, adjust the \texttt{cycles} line in \texttt{dosbox-staging.conf}:

\begin{tomlcode}
[cpu]
cycles = fixed 8000   # Increase if too slow, decrease if too fast
\end{tomlcode}

The original Hugo required roughly a 386SX at 25\,MHz. Start with \texttt{cycles = fixed 8000} and calibrate.

\section{Step 6: Exit cleanly}

Two \texttt{Esc} presses bring up Hugo's menu. The \emph{Quit} option returns to the DOSBox prompt; the \texttt{exit} command in the \texttt{[autoexec]} block closes DOSBox automatically.

For an emergency exit (for example if the game hangs): \texttt{Ctrl}+\texttt{F9} is the \emph{kill switch} for DOSBox-Staging.

\section{Troubleshooting}

\subsection{Hugo hangs on the intro or the cutscenes}

DOSBox-Staging on ARM occasionally has problems decoding VGA animations. Pressing any key to skip the intro usually solves it. If it persists, lower the output resolution:

\begin{tomlcode}
[sdl]
windowresolution = 640x480
\end{tomlcode}

\subsection{The audio stutters}

Increase the mixer buffer:

\begin{tomlcode}
[mixer]
prebuffer = 50
blocksize = 2048
\end{tomlcode}

If stuttering remains, lower the quality: \texttt{rate = 22050}.

\subsection{Black screen at startup}

The KMSDRM driver is probably not available in the installed SDL2 version. Alternatives:

\begin{shellcode}
SDL_VIDEODRIVER=fbcon dosbox-staging
\end{shellcode}

Or install the minimal desktop:

\begin{shellcode}
sudo apt install -y xinit openbox
\end{shellcode}

And launch DOSBox inside \texttt{startx}.

\subsection{Hugo starts but doesn't respond to the keyboard}

\begin{itemize}
  \item Verify that \texttt{core = dynamic} is present in the \texttt{[cpu]} section.
  \item Some titles require \texttt{cputype = 386} instead of \texttt{486\_slow}.
\end{itemize}

\section{Phase wrap-up}

At this point the game runs with a USB keyboard. The next phase replaces that USB keyboard with a virtual one controlled from the phone.
